Il progetto FitTag è un ecosistema IoT progettato per il monitoraggio intelligente delle attività sportive.
Il progetto è composto da due macro-sistemi strettamente interconnessi: un bracciale hardware intelligente (FitTag Band) e un'applicazione mobile dedicata per l'interazione con l'utente (FitTag App).

\paragraph{FitTag Band}
Il dispositivo wearable è un sistema indipendente dotato di batteria, unità di misuramento inerziale e pulsossimetro.
La principale funzionalità del sistema è la registrazione di sessioni di allenamento, arricchita dal riconoscimento automatico dell'attività sportiva che l'utente sta svolgendo.
Questa capacità è resa possibile dall'implementazione on-device di un modello di Machine Learning ottimizzato per l'inferenza in tempo reale.

\paragraph{FitTag App}
L'applicazione mobile costituisce l'interfaccia utente dell'ecosistema, progettata per operare in stretta sinergia con il bracciale tramite comunicazione wireless Bluetooth Low Energy (BLE).
La principale funzionalità del software è la gestione remota del dispositivo e delle sessioni di allenamento, affiancata dalla raccolta e visualizzazione delle metriche sportive.
Questa interazione bidirezionale permette all'app sia di inviare comandi operativi al wearable, sia di restituire all'utente un quadro immediato dei parametri vitali e delle attività predette dal modello di Machine Learning.